\section{其他拟合选项与稳定性}\label{sec:highorder}

在唯象分析中，能得到的信息量与精度显然取决于格点数据的质量。就现有数据而言，我们希望在拟合中纳入更多物理，直至分析不再给出有用信息为止，以此逼近分析的极限。

\begin{figure}[tbp]
\centering
\includegraphics[width=9cm]{images/chisqmap_Eq26_clover_Pion.pdf}
\includegraphics[width=8cm]{images/Eq26_clover_Pion_Expgzmm0zdd.pdf}
\caption{
用式~(\ref{eq:logM_2}) 重正化 clover 作用量、无 HYP 涂抹时 $\pi$ 态中的 $O_{\gamma_t}(z)$ 矩阵元。上图：$\chi^2$ 随 $\lambda$ 与 $\Lambda_{\rm QCD}$ 的分布图，$\langle  \chi^{2}/{\rm d.o.f.} \rangle_{z}$ 是各 $z$ 拟合的 $\chi^{2}/{\rm d.o.f.}$ 的平均值。下图：沿“小卡方带”若干 ($\Lambda_{\rm QCD}$, $\lambda$) 参数组的重正化矩阵元。
}
\label{fig:Re_overlap_quark_2}
\end{figure}

\begin{figure}[tbp]
\centering
\includegraphics[width=9cm]{images/chisqmap_Eq27_clover_Pion.pdf}
\includegraphics[width=8cm]{images/Eq27_clover_Pion_Expgzmm0zdd.pdf}
\caption{同图~\ref{fig:Re_overlap_quark_2}，但用式 (\ref{eq:logM_3}) 拟合。
}
\label{fig:Re_overlap_quark_3}
\end{figure}

我们想回答的第一个问题是：能否获得关于线性发散的更精确的信息？迄今我们使用的是单圈微扰 QCD 预言，而把 $\Lambda_{\rm QCD}$ 作为自由参数以部分计入高阶修正。不过也可以考虑直接至少计入二阶修正。这里的策略是固定 $k$ 为其微扰值，同时变动 $\Lambda_{\rm QCD}$ 以及由新参数 $\lambda$ 刻画的显式二阶修正
\begin{align}\label{eq:logM_2}
\ln \mathcal{M}(z,a) = \frac{k' z}{a} \tilde{\alpha_{s}}(1+\lambda \tilde{\alpha_{s}})
+ g(z) + ...
\end{align}
其中省略的各项与式~(\ref{eq:logM}) 中先前的相同。$k'$ 通过 $-2\pi/b_0$ 与 $k$ 相联系（因为 $-k'2\pi/b_0=k$）。另一方面，可以用参数 $\lambda$ 纳入某些高阶修正~\cite{Bauer:2011ws},

\begin{align}\label{eq:logM_3}
\ln \mathcal{M}(z,a) = \frac{k' z}{a} \frac{\tilde{\alpha_{s}}}{1-\lambda \tilde{\alpha_{s}}}
+ g(z) + ...
\end{align}
为保持一致，现在须用到精确到二阶 $\beta$ 函数 $b_{1}=102-\frac{38}{3} n_{f}$ 的 QCD 跑动耦合常数 $\tilde{\alpha_{s}}$，
\bea\label{eq:alpha_s2}
\tilde{\alpha_{s}}=\frac{4 \pi}{b_0 \ln (\frac{1}{a^{2} \Lambda_{\rm QCD}^{2}})}\left(1-\frac{b_1}{b_0^{2}} \times \frac{\ln[\ln (\frac{1}{a^{2} \Lambda_{\rm QCD}^{2}})]}{\ln(\frac{1}{a^{2} \Lambda_{\rm QCD}^{2}})}\right).
\eea
因此，拟合中又有两个全局参数 $\lambda$ 与 $\Lambda_{\rm QCD}$。

采用上述参数化的拟合过程与第~\ref{sec:reofpro} 节所述类似，结果示于图~\ref{fig:Re_overlap_quark_2} 和图~\ref{fig:Re_overlap_quark_3}。
可以看出，$\lambda$ 与 $\Lambda_{\rm QCD}$ 的允许范围相当大，且二者强烈关联。选出若干相互关联的参数组并匹配到微扰结果后，所得残差矩阵元示于下图；不同选择之间几乎看不出差别。而且这些结果与不含额外 $\lambda$ 参数的拟合（图~\ref{fig:Re_overlap_quark_dff}）基本相同。由此我们得出结论：就现有数据而言，我们的结果对高阶修正是稳定的；只有更精确的数据才可能分辨出高阶拟合的差异。

\begin{figure}[tbp]
\centering
\includegraphics[width=9cm]{images/chisqmap_Eq29_clover_Pion.pdf}
\includegraphics[width=8cm]{images/Eq29_clover_Pion_Expgzmm0zdd.pdf}
\caption{
用式~(\ref{eq:logM_4}) 重正化 clover 作用量、无 HYP 涂抹时 $\pi$ 态中的 $O_{\gamma_t}(z)$ 矩阵元。上图：$\chi^2$ 随 $\Lambda_{\rm QCD1}$ 与 $\Lambda_{\rm QCD2}$ 的分布图，$\langle  \chi^{2}/{\rm d.o.f.} \rangle_{z}$ 是各 $z$ 拟合的 $\chi^{2}/{\rm d.o.f.}$ 的平均值。下图：沿“小卡方带”若干 ($\Lambda_{\rm QCD1}$, $\Lambda_{\rm QCD2}$) 参数组的重正化矩阵元。
}
\label{fig:Re_overlap_quark_4}
\end{figure}

\begin{figure}[tbp]
\centering
\includegraphics[width=8cm]{images/clover_Pion_Expgzmm0zdd_dm.pdf}
\caption{
clover 作用量、无 HYP 涂抹时 $\pi$ 态基于不同拟合函数得到的重正化 $O_{\gamma_t}(z)$ 关联：式~(\ref{eq:logM})（紫）、式~(\ref{eq:logM_2})（红）、式~(\ref{eq:logM_3})（蓝）、式~(\ref{eq:logM_4})（绿）。
}
\label{fig:Re_overlap_quark_dff}
\end{figure}

另一个可探索的可能性是：由 $\Lambda_{\rm QCD}$ 代表的高阶修正对不同系综可能不同。因此可以对 MILC 与 RBC 数据讨论不同的 $\Lambda_{\rm QCD}$ 取值，
\begin{align}\label{eq:logM_4}
\ln \mathcal{M}(z,a) = \frac{k z}{a \ln[a \Lambda_{\rm QCD1,2}]} + g(z) + f_{1,2} (z) a \nonumber\\
+\frac{3 C_{F}}{b_0} \ln [\frac{\ln [1 /(a \Lambda_{\rm QCD1,2})]}{\ln [\mu / \Lambda_{\rm QCD1,2}]}] + \ln [1+\frac{d}{\ln (a \Lambda_{\rm QCD1,2})}],
\end{align}
其中固定参数 $k$，得到双参数拟合 $\Lambda_{\rm QCD1}$ 与 $\Lambda_{\rm QCD2}$。
使用上式的结果示于图~\ref{fig:Re_overlap_quark_4}。如图所示，两个参数强烈关联且互相成正比，二者的差被限制在一个很小的范围内；而且取不同的关联参数组 ($\Lambda_{\rm QCD1}$, $\Lambda_{\rm QCD2}$)，最终的重正化结果不变。

在图~\ref{fig:Re_overlap_quark_dff} 中，我们展示了基于不同拟合函数得到的重正化矩阵元结果。所有这些拟合函数都给出同一结果。因此，就现有数据而言，没有必要在拟合函数中考虑其他高阶项，式~(\ref{eq:logM}) 已经足够好。这也支持第~\ref{ReUncer} 节的论断：可以用领头阶拟合得到的 $\Lambda_{\rm QCD}$ 部分地代表高阶修正。

最后，我们考虑拟合公式的 $O(a^{2})$ 修正。这些拟合并不稳定，说明对当前的数据质量而言拟合参数引入得过多。结论是：需要更高精度的数据才能约束 $O(a^2)$ 项。也可以在线性修正之外改用 $O(a^{2})$ 修正来做拟合；对我们研究的情形而言，这样做的物理动机较弱。
